\documentclass[12pt]{article}

% --- Packages ---
\usepackage[margin=1in]{geometry}
\usepackage{amsmath,amssymb,amsthm}
\usepackage{mathtools}
\usepackage{graphicx}
\usepackage{booktabs}
\usepackage{natbib}
\usepackage{hyperref}
\usepackage{cleveref}
\usepackage{enumitem}
\usepackage{multirow}
\usepackage{array}

% --- Theorem environments ---
\newtheorem{theorem}{Theorem}
\newtheorem{lemma}[theorem]{Lemma}
\newtheorem{proposition}[theorem]{Proposition}
\newtheorem{corollary}[theorem]{Corollary}
\newtheorem{conjecture}[theorem]{Conjecture}
\theoremstyle{definition}
\newtheorem{definition}[theorem]{Definition}
\newtheorem{example}[theorem]{Example}
\theoremstyle{remark}
\newtheorem{remark}[theorem]{Remark}

% --- Notation shortcuts ---
\newcommand{\E}{\mathbb{E}}
\newcommand{\Var}{\mathrm{Var}}
\newcommand{\tr}{\mathrm{tr}}
\newcommand{\diag}{\mathrm{diag}}
\newcommand{\rank}{\mathrm{rank}}
\newcommand{\Rset}{\mathbb{R}}
\newcommand{\bone}{\mathbf{1}}
\newcommand{\blam}{\boldsymbol{\lambda}}
\newcommand{\btheta}{\boldsymbol{\theta}}
\DeclareMathOperator*{\argmax}{arg\,max}
\DeclareMathOperator*{\argmin}{arg\,min}
\DeclareMathOperator{\sgn}{sgn}

% --- Title ---
\title{Information Geometry of Coherent Systems:\\
       Eigenvalue Structure and Component Estimability}
\author{Alexander Towell\thanks{Department of Computer Science,
  Southern Illinois University Edwardsville.
  Email: \texttt{lex@metafunctor.com}.
  ORCID: 0000-0001-6443-9897.}}
\date{}

\begin{document}
\maketitle

\begin{abstract}
We study the Fisher information matrix $\mathcal{I}(\blam)$ for
component rate estimation from system-level lifetime data (Scheme~0)
across a range of coherent system structures.  Our main contributions
are threefold.  First, we prove that for series systems with
exponential components, $\rank(\mathcal{I}) = 1$---only the total
failure rate is identifiable---while for $k$-out-of-$n$ systems with
$1 \leq k < n$ and distinct rates, $\rank(\mathcal{I}) = n$, so all
component rates are identifiable.  Second, we show that the eigenvalue
structure of $\mathcal{I}$ provides a sharper characterization of
estimation difficulty than the determinant: the leading eigenvalue
corresponds to the total-rate direction (always well-estimated), while
trailing eigenvalues capture rate-difference directions (poorly
estimated for some structures).  The condition number
$\kappa(\mathcal{I}) = \lambda_{\max}/\lambda_{\min}$ predicts
mean RMSE with $R^2 > 0.9$ on a log-log scale, whereas the
determinant---the basis of D-optimality---is misleading when one
eigenvalue dominates.  Third, we establish that component-symmetric
systems (those whose structure function is invariant under component
permutation) exhibit uniformly better conditioning of $\mathcal{I}$
than their asymmetric counterparts, and we identify ``information
deserts'' in asymmetric structures where edge components are
effectively unestimable.  A simulation study across eleven system
structures with three and five components provides numerical
confirmation.
\end{abstract}

\noindent\textbf{Keywords:} coherent systems, Fisher information,
eigenvalue structure, condition number, identifiability, $k$-out-of-$n$
systems, consecutive-$k$ systems, component estimation, information
geometry, optimal design

%%% ====================================================================
\section{Introduction}
\label{sec:introduction}

The problem of estimating component lifetime distributions from
system-level observations has attracted sustained attention in
reliability theory.  When only the system lifetime $T$ is observed---a
setting we call \emph{Scheme~0}---the analyst faces an information
bottleneck: the system structure compresses the $m$-dimensional
component lifetime vector into a single observable.  How much
information about the individual component parameters survives this
compression depends entirely on the system structure.

For series systems ($T = \min(T_1, \ldots, T_m)$), the answer is stark.
The foundational result of \citet{Tsiatis1975} shows that, under
independence, the component distributions are not identifiable from
the distribution of $T$ alone: infinitely many joint distributions
produce the same marginal for the minimum.  For exponential components,
only the total rate $\sum_j \lambda_j$ is determined by the system
lifetime distribution.  In terms of Fisher information, this means
$\rank(\mathcal{I}) = 1$---the information matrix is rank-deficient,
and $m - 1$ parameter directions carry zero information.

For parallel systems ($T = \max(T_1, \ldots, T_m)$), the situation
is qualitatively different.  \citet{BasuGhosh1980} showed that the
distribution of the maximum, under independence, uniquely determines
the component distributions.  More generally, \citet{Nowik1990}
characterized which coherent system structures preserve
identifiability.  But identifiability---the assertion that the Fisher
information matrix has full rank---says nothing about the magnitude
or distribution of information across parameter directions.  A system
can be identifiable yet have a condition number so large that some
component rates are effectively unestimable at practical sample sizes.

This paper develops a systematic theory of estimation difficulty for
coherent systems with exponential components under Scheme~0.  Our
perspective is information-geometric: instead of asking merely
whether the Fisher information matrix is positive definite, we study
its eigenvalue structure---the spectrum, condition number, and
eigenvector directions---as a function of system topology.

\paragraph{Contributions.}
The specific contributions are:
\begin{enumerate}[leftmargin=2em]
  \item \emph{Rank characterization} (Theorems~\ref{thm:series-rank}
    and~\ref{thm:kofn-rank}):
    We prove that series systems have $\rank(\mathcal{I}) = 1$ and
    that $k$-out-of-$n$ systems with $1 \leq k < n$ and distinct
    rates have $\rank(\mathcal{I}) = n$.

  \item \emph{Eigenvalue structure} (\S\ref{sec:eigenvalue}):
    We show that the leading eigenvalue of $\mathcal{I}$ corresponds
    to the total-rate direction $\bone/\sqrt{m}$, while trailing
    eigenvalues capture contrast directions.  The decay rate of the
    eigenvalue spectrum depends on the system signature.

  \item \emph{Condition number as universal predictor}
    (\S\ref{sec:det-vs-kappa}):
    We demonstrate numerically that $\log(\text{RMSE})$ is linearly
    predicted by $\log(\kappa(\mathcal{I}))$ with $R^2 > 0.9$ across
    eleven system structures, while the determinant (D-optimality) can
    be misleading.

  \item \emph{Component symmetry and information distribution}
    (\S\ref{sec:symmetry}):
    We establish that component-symmetric structures (those invariant
    under permutation of component labels) have better-conditioned
    information matrices than asymmetric structures, and we
    characterize the ``information deserts'' that arise in
    consecutive-$k$ systems.

  \item \emph{Numerical study} (\S\ref{sec:simulations}):
    A comprehensive simulation across eleven coherent structures
    (series, parallel, $k$-out-of-$n$, bridge, consecutive-$k$)
    provides numerical support for the theoretical results and reveals
    the quantitative relationship between eigenvalue structure and
    estimation quality.
\end{enumerate}

\paragraph{Relation to prior work.}
The identifiability of component distributions from system data has been
studied by \citet{BasuGhosh1980}, \citet{Meilijson1981}, and
\citet{Nowik1990}, among others.  Estimation from system data has been
pursued via signature-based methods
\citep{BhattacharyaSamaniego2010, NgNavarroBalakrishnan2017,
YangNgBalakrishnan2019, HallJinSamaniego2015} and parametric methods
\citep{LinUsherGuess1993, SarhanElBassiouny2003, BheringPereiraPolpo2014}.
\citet{CramerNavarro2015} established a connection between coherent
systems and progressive Type-II censoring.  The present paper differs
from all of these in its focus: we study the \emph{geometry} of the
information loss, not just its presence or absence, and we show that
the eigenvalue structure of the Fisher information matrix provides a
complete characterization of practical estimability.

The connection to optimal experimental design
\citep{Kiefer1974, Pukelsheim2006} arises naturally: the system
structure plays the role of the design, and the Fisher information
matrix is the design matrix.  Our finding that D-optimality (based on
$\det(\mathcal{I})$) can be misleading for system selection parallels
known limitations of D-optimality in experimental design when the
eigenvalue spectrum is highly skewed \citep{Atkinson2007}.

\paragraph{Organization.}
Section~\ref{sec:prelim} reviews coherent systems and fixes notation.
Section~\ref{sec:fisher} derives the Fisher information matrix for
general coherent systems with exponential components and proves the
rank theorems.
Section~\ref{sec:eigenvalue} analyzes the eigenvalue structure.
Section~\ref{sec:det-vs-kappa} compares the determinant and condition
number as measures of estimation difficulty.
Section~\ref{sec:symmetry} studies the role of component symmetry.
Section~\ref{sec:simulations} presents the simulation study.
Section~\ref{sec:discussion} discusses connections, extensions, and
open problems.
Proofs of selected results appear in the appendix.


%%% ====================================================================
\section{Preliminaries}
\label{sec:prelim}

We review the basics of coherent systems, following
\citet{BarlowProschan1975} and \citet{Samaniego2007}.

\subsection{Coherent systems}
\label{sec:coherent-def}

\begin{definition}[Structure function]
\label{def:structure-function}
A \emph{structure function} is a mapping
$\phi\colon \{0,1\}^m \to \{0,1\}$ where $\phi(\mathbf{x})$ indicates
whether the system functions ($\phi = 1$) or fails ($\phi = 0$) when
the component state vector is $\mathbf{x} = (x_1, \ldots, x_m)$ with
$x_j = 1$ meaning component~$j$ is operational and $x_j = 0$ meaning
it has failed.
\end{definition}

\begin{definition}[Coherent system]
\label{def:coherent}
A structure function $\phi$ is \emph{coherent} if:
\begin{enumerate}[label=(\roman*)]
  \item $\phi$ is nondecreasing in each argument (improving a component
    cannot harm the system);
  \item every component is relevant: for each $j$, there exists
    $\mathbf{x}_{-j}$ such that
    $\phi(1_j, \mathbf{x}_{-j}) \neq \phi(0_j, \mathbf{x}_{-j})$.
\end{enumerate}
\end{definition}

\begin{definition}[Path and cut sets]
\label{def:path-cut}
A \emph{path set} $P \subseteq \{1, \ldots, m\}$ is a set of components
whose simultaneous functioning guarantees system functioning:
$\phi(\mathbf{x}) = 1$ whenever $x_j = 1$ for all $j \in P$.
A path set is \emph{minimal} if no proper subset is also a path set.
Dually, a \emph{cut set} $K$ is a set of components whose simultaneous
failure guarantees system failure, and a minimal cut set has no proper
subset that is also a cut set.
\end{definition}

The system functions if and only if at least one minimal path set is
fully operational; the system fails if and only if at least one minimal
cut set has fully failed.  If $\mathcal{P} = \{P_1, \ldots, P_p\}$
denotes the collection of minimal path sets and
$\mathcal{K} = \{K_1, \ldots, K_q\}$ the minimal cut sets, then
\begin{equation}\label{eq:lifetime-paths-cuts}
  T = \max_{P \in \mathcal{P}} \min_{j \in P} T_j
    = \min_{K \in \mathcal{K}} \max_{j \in K} T_j.
\end{equation}

\begin{definition}[Critical state]
\label{def:critical}
Component~$j$ is \emph{critical} (or \emph{pivotal}) in state
$\mathbf{x}_{-j}$ if
$\phi(1_j, \mathbf{x}_{-j}) = 1$ and $\phi(0_j, \mathbf{x}_{-j}) = 0$.
We write $\mathcal{C}_j$ for the set of states $\mathbf{x}_{-j}$ in
which $j$ is critical.
\end{definition}

\begin{example}[Standard structures]\label{ex:structures}
\begin{enumerate}[label=(\alph*)]
  \item \emph{Series}: $\phi(\mathbf{x}) = \prod_j x_j$.
    Single path set $\{1, \ldots, m\}$; $m$ singleton cut sets.
  \item \emph{Parallel}: $\phi(\mathbf{x}) = 1 - \prod_j(1 - x_j)$.
    $m$ singleton path sets; single cut set $\{1, \ldots, m\}$.
  \item \emph{$k$-out-of-$n$}: system works iff at least $k$ components
    work.  Path sets are all $k$-subsets; cut sets are all
    $(n - k + 1)$-subsets.  Series is $n$-out-of-$n$; parallel is
    $1$-out-of-$n$.
  \item \emph{Consecutive-$k$-of-$n$}: system fails if $k$ consecutive
    components all fail \citep{ChiangNiu1981}.  Cut sets are the
    $(n - k + 1)$ runs of $k$ consecutive indices.
  \item \emph{Bridge} (5~components): minimal path sets
    $\{1,4\}$, $\{2,5\}$, $\{1,3,5\}$, $\{2,3,4\}$.  A classical
    non-series-parallel structure.
\end{enumerate}
\end{example}

\subsection{System signature}
\label{sec:signature}

\begin{definition}[Signature {\citep{Samaniego1985}}]
\label{def:signature}
The \emph{signature} of a coherent system with $m$ components is the
vector $\mathbf{s} = (s_1, \ldots, s_m)$ where
$s_i = P(T = T_{(i)})$ under the assumption that the component
lifetimes $T_1, \ldots, T_m$ are i.i.d.\ continuous, and $T_{(i)}$
denotes the $i$-th order statistic.  The signature depends only on
the structure function, not on the component distributions.
\end{definition}

The signature encodes where in the failure sequence the system fails.
For a series system, $\mathbf{s} = (1, 0, \ldots, 0)$: the system
fails at the first failure.  For a parallel system,
$\mathbf{s} = (0, \ldots, 0, 1)$: the system fails at the last
failure.  For a $k$-out-of-$n$ system,
$s_i = \mathbf{1}(i = n - k + 1)$.

\subsection{Component lifetime model}
\label{sec:component-model}

Throughout this paper, we consider $m$ independent components with
exponential lifetimes:
\begin{equation}\label{eq:component-exp}
  T_j \sim \mathrm{Exp}(\lambda_j), \qquad
  F_j(t) = 1 - e^{-\lambda_j t}, \qquad
  f_j(t) = \lambda_j e^{-\lambda_j t}, \qquad
  S_j(t) = e^{-\lambda_j t},
\end{equation}
for $j = 1, \ldots, m$ and $t > 0$.  The parameter vector is
$\blam = (\lambda_1, \ldots, \lambda_m) \in \Rset_{>0}^m$.
Throughout, we order the rates as
$\lambda_1 \leq \lambda_2 \leq \cdots \leq \lambda_m$
(smallest rate = most reliable component = largest mean lifetime).

\subsection{Component symmetry}
\label{sec:symmetry-def}

\begin{definition}[Component-symmetric system]
\label{def:component-symmetric}
A coherent system is \emph{component-symmetric} if its structure
function is invariant under every permutation of component labels:
$\phi(\mathbf{x}_{\sigma}) = \phi(\mathbf{x})$ for all permutations
$\sigma \in \mathfrak{S}_m$ and all $\mathbf{x} \in \{0,1\}^m$.
\end{definition}

The $k$-out-of-$n$ systems are the canonical component-symmetric
structures: whether the system works depends only on how many
components work, not which ones.  Consecutive-$k$ systems are
\emph{not} component-symmetric because interior and edge components
play structurally different roles.


%%% ====================================================================
\section{Fisher Information for Coherent Systems}
\label{sec:fisher}

\subsection{System density via critical-state decomposition}
\label{sec:system-density}

For a coherent system with structure function $\phi$ and independent
components with distributions $(f_j, F_j, S_j)$, the system density
admits the \emph{critical-state decomposition}:
\begin{equation}\label{eq:f-sys-general}
  f_T(t; \blam) = \sum_{j=1}^m f_j(t) \sum_{\mathbf{x}_{-j} \in \mathcal{C}_j}
  \prod_{\substack{i \neq j \\ x_i = 1}} S_i(t)
  \prod_{\substack{i \neq j \\ x_i = 0}} F_i(t),
\end{equation}
where $\mathcal{C}_j$ is the set of critical states for component~$j$
(Definition~\ref{def:critical}).  This representation expresses the
system density as a sum over ``who failed last'' (index $j$) and ``what
was the system state when $j$ failed'' (state $\mathbf{x}_{-j}$).

For exponential components, each factor in \eqref{eq:f-sys-general}
is either $e^{-\lambda_i t}$ (survival) or $1 - e^{-\lambda_i t}$
(CDF), and the density is a linear combination of terms of the form
$\lambda_j \, e^{-\alpha t}$ for various rate sums $\alpha$.

\subsection{Log-likelihood under Scheme~0}
\label{sec:loglik}

Under Scheme~0, we observe $n$ independent copies
$t_1, \ldots, t_n$ of the system lifetime.  The log-likelihood is
\begin{equation}\label{eq:loglik}
  \ell(\blam) = \sum_{i=1}^n \log f_T(t_i; \blam).
\end{equation}

\begin{definition}[Observed Fisher information]
\label{def:fisher-obs}
The \emph{observed Fisher information matrix} at $\hat{\blam}$ is
\begin{equation}\label{eq:fisher-obs}
  \hat{\mathcal{I}}(\hat{\blam})
  = -\frac{\partial^2 \ell}{\partial \blam \,\partial \blam^\top}
  \bigg|_{\blam = \hat{\blam}}.
\end{equation}
Under standard regularity conditions,
$\hat{\mathcal{I}} / n \to \mathcal{I}(\blam_0)$ in probability, where
$\mathcal{I}(\blam) = \E_{\blam}\bigl[-\partial^2 \log f_T(T; \blam) /
\partial \blam \,\partial \blam^\top\bigr]$
is the expected Fisher information per observation.
\end{definition}

We write $\mathcal{I}$ for the (expected or observed, as context
dictates) Fisher information matrix, suppressing the argument $\blam$
when clear.

\subsection{Rank of the Fisher information: Series systems}
\label{sec:series-rank}

\begin{theorem}[Series rank deficiency]
\label{thm:series-rank}
For a series system of $m$ independent exponential components with
rates $\lambda_1, \ldots, \lambda_m > 0$, observed under Scheme~0,
\begin{equation}\label{eq:series-rank}
  \rank(\mathcal{I}(\blam)) = 1.
\end{equation}
The unique nonzero eigenvalue corresponds to the eigenvector
$\bone / \sqrt{m}$, the total-rate direction.
\end{theorem}

\begin{proof}
For a series system, $T = \min(T_1, \ldots, T_m)$, and by independence,
$T \sim \mathrm{Exp}(\lambda_+)$ where
$\lambda_+ = \sum_{j=1}^m \lambda_j$.
The system density is $f_T(t) = \lambda_+ e^{-\lambda_+ t}$, and
the log-likelihood contribution from a single observation $t$ is
\[
  \log f_T(t) = \log \lambda_+ - \lambda_+ t.
\]
Now $\lambda_+$ depends on $\blam$ only through the linear functional
$\blam \mapsto \bone^\top \blam$, so
\[
  \frac{\partial \log f_T}{\partial \lambda_j}
  = \frac{1}{\lambda_+} - t
\]
for every $j$.  The score vector is
$\nabla_{\blam} \log f_T = (1/\lambda_+ - t)\,\bone$, which lies in
the one-dimensional subspace spanned by $\bone$.

The Fisher information matrix is
\[
  \mathcal{I}(\blam) = \E\bigl[(\nabla \log f_T)(\nabla \log f_T)^\top\bigr]
  = \frac{1}{\lambda_+^2}\,\bone\bone^\top,
\]
using $\Var(T) = 1/\lambda_+^2$.  This is a rank-1 matrix with
eigenvalue $m/\lambda_+^2$ (with eigenvector $\bone/\sqrt{m}$) and
$m-1$ zero eigenvalues.

The rank-1 structure reflects the fundamental non-identifiability
result of \citet{Tsiatis1975}: the system data determines only
$\lambda_+$, and no individual rate can be separated from the total.
\end{proof}

\begin{remark}
The complete-data Fisher information for $n$ i.i.d.\
$\mathrm{Exp}(\lambda_j)$ observations is
$\mathcal{I}_j^{\mathrm{comp}} = n/\lambda_j^2$, independent across
components.  The series system projects the full $m$-dimensional
parameter space onto a one-dimensional subspace, destroying
$(m-1)/m$ of the eigenvalue directions.
\end{remark}

\subsection{Rank of the Fisher information: $k$-out-of-$n$ systems}
\label{sec:kofn-rank}

\begin{theorem}[Full rank for $k$-out-of-$n$]
\label{thm:kofn-rank}
For a $k$-out-of-$n$ system with $1 \leq k < n$ and independent
exponential components with distinct rates
$0 < \lambda_1 < \lambda_2 < \cdots < \lambda_n$, observed under
Scheme~0,
\begin{equation}\label{eq:kofn-rank}
  \rank(\mathcal{I}(\blam)) = n.
\end{equation}
All component rates are identifiable from the system lifetime
distribution.
\end{theorem}

\begin{proof}
We establish identifiability by showing that the system density
$f_T(t; \blam)$ uniquely determines $\blam$, which implies that
$\mathcal{I}(\blam)$ is positive definite at interior points of the
parameter space.

For a $k$-out-of-$n$ system, $T$ is the $(n - k + 1)$-th order
statistic of $(T_1, \ldots, T_n)$.  By independence, the system CDF is
\begin{equation}\label{eq:kofn-cdf}
  F_T(t) = \sum_{i=n-k+1}^{n} \sum_{|S|=i} \prod_{j \in S} F_j(t)
  \prod_{j \notin S} S_j(t).
\end{equation}
For exponential components, each product in~\eqref{eq:kofn-cdf} is of
the form $e^{-\alpha t}$ multiplied by a polynomial in terms
$e^{-\lambda_j t}$.  After expansion, $F_T(t)$ is a finite linear
combination of exponential functions $e^{-\beta t}$ where each $\beta$
is a sum of a subset of the $\lambda_j$'s.

The key observation is that for $k < n$, the system CDF involves
products that mix survival functions $S_j = e^{-\lambda_j t}$ with
CDFs $F_j = 1 - e^{-\lambda_j t}$, producing a rich collection of
rate sums.  When the $\lambda_j$ are distinct, the subset sums
$\beta_S = \sum_{j \in S} \lambda_j$ are generically distinct (i.e.,
for all $\blam$ outside a measure-zero set), so $F_T(t)$ is a linear
combination of exponentials with distinct rates.

A finite linear combination of exponentials with distinct rates is
uniquely determined by the rates and coefficients (this follows from
the linear independence of $\{e^{-\beta t} : \beta > 0\}$ over the
reals).  Since the coefficients in the expansion of $F_T$ are
determined by $\blam$ and the $\beta_S$ are distinct sums of
$\lambda_j$'s, the mapping $\blam \mapsto F_T(\cdot; \blam)$ is
injective.

The case $k = 1$ (parallel system) was established by
\citet{BasuGhosh1980}, who showed that the product
$\prod_j (1 - e^{-\lambda_j t})$ has a unique factorization when
the rates are distinct.  The argument extends to general $k < n$ by
the same linear-independence reasoning.

Since identifiability holds generically, the Fisher information matrix
$\mathcal{I}(\blam)$ has full rank at every $\blam$ with distinct
components.
\end{proof}

\begin{remark}
The case $k = n$ (series) is excluded: in a series system,
$F_T(t) = 1 - \prod_j S_j(t) = 1 - e^{-\lambda_+ t}$, which
depends on $\blam$ only through $\lambda_+$, confirming the rank-1
result of Theorem~\ref{thm:series-rank}.
\end{remark}

\begin{corollary}\label{cor:general-coherent}
For any coherent system that is not a series system, if the structure
function involves at least one minimal path set of size strictly less
than $m$, then generically $\rank(\mathcal{I}(\blam)) = m$ for
exponential components with distinct rates.
\end{corollary}

\begin{proof}[Proof sketch]
A minimal path set of size $r < m$ means that the system can function
with $m - r$ components failed.  In the critical-state decomposition
\eqref{eq:f-sys-general}, this creates terms where some components
contribute $F_i(t)$ (failed) while others contribute $S_i(t)$
(survived), generating a mixture of exponentials with different
rate sums.  The same linear-independence argument as in
Theorem~\ref{thm:kofn-rank} gives injectivity generically.

The only coherent system where all path sets have size $m$ is the
series system (the single path set is the entire component set).
\end{proof}


%%% ====================================================================
\section{Eigenvalue Structure of the Fisher Information}
\label{sec:eigenvalue}

Having established \emph{when} the Fisher information has full rank
(Theorems~\ref{thm:series-rank}--\ref{thm:kofn-rank}), we now study
\emph{how} the information is distributed across parameter directions.
The eigenvalue decomposition
$\mathcal{I} = \sum_{k=1}^m \mu_k \mathbf{v}_k \mathbf{v}_k^\top$
(with $\mu_1 \geq \mu_2 \geq \cdots \geq \mu_m \geq 0$) reveals
which linear combinations of the $\lambda_j$ are well-determined and
which are poorly determined by the data.

\subsection{The total-rate direction}
\label{sec:total-rate}

\begin{proposition}[Leading eigenvector]
\label{prop:leading-ev}
For any coherent system with independent exponential components under
Scheme~0, the leading eigenvector $\mathbf{v}_1$ of
$\mathcal{I}(\blam)$ is approximately aligned with the total-rate
direction $\bone / \sqrt{m}$.  More precisely, if $\mu_1$ denotes
the largest eigenvalue, then
\begin{equation}\label{eq:leading-ev-bound}
  |\mathbf{v}_1^\top \bone / \sqrt{m}| \geq 1 - O(\mu_2 / \mu_1),
\end{equation}
where $\mu_2$ is the second-largest eigenvalue.
\end{proposition}

\begin{proof}
The score vector $\nabla_{\blam} \log f_T(t; \blam)$ has the
structure
\[
  \frac{\partial}{\partial \lambda_j} \log f_T(t; \blam)
  = \frac{1}{f_T} \frac{\partial f_T}{\partial \lambda_j}.
\]
For any coherent system, the derivative $\partial f_T / \partial \lambda_j$
involves the density of component~$j$ weighted by the failure/survival
probabilities of the other components.  In the exponential case, the
dominant contribution at large $t$ comes from the slowest-decaying
exponential terms, which involve $\lambda_+$ (the total rate).
Consequently, $\E[(\partial \log f_T / \partial \lambda_j)^2]$ is
dominated by a component proportional to $1/\lambda_+^2$, common to
all $j$.

The Fisher information matrix therefore decomposes as
\[
  \mathcal{I} = \alpha\, \bone\bone^\top + \mathbf{R},
\]
where $\alpha > 0$ is a scalar and $\mathbf{R}$ is a residual matrix
capturing the system-specific information about rate differences.
The eigenvalue $\mu_1 \approx m\alpha + \lambda_{\max}(\mathbf{R})$
is dominated by the $\alpha\, \bone\bone^\top$ term, so $\mathbf{v}_1$
is close to $\bone/\sqrt{m}$.  The alignment improves as $\mu_1 / \mu_2$
increases.
\end{proof}

\begin{remark}
Proposition~\ref{prop:leading-ev} explains a universal feature of
our numerical experiments: across all eleven system structures tested,
the leading eigenvector has inner product $> 0.99$ with
$\bone / \sqrt{m}$.  The total rate $\lambda_+$ is always
well-estimated; the difficulty lies entirely in the trailing
eigenvalue directions.
\end{remark}

\subsection{Trailing eigenvalues and contrast directions}
\label{sec:trailing-ev}

The trailing eigenvectors $\mathbf{v}_2, \ldots, \mathbf{v}_m$---the
contrast directions---capture the information about \emph{differences}
between component rates.  These are orthogonal to the total-rate
direction and correspond to the parameter combinations that are
hardest to estimate.

\begin{proposition}[Eigenvalue decay and signature]
\label{prop:ev-decay}
For a $k$-out-of-$n$ system with i.i.d.\ component rates, the
eigenvalue ratio $\mu_{\min} / \mu_{\max}$ depends on the signature
vector $\mathbf{s}$.  Specifically:
\begin{enumerate}[label=(\roman*)]
  \item For the parallel system ($k = 1$, signature
    $\mathbf{s} = (0, \ldots, 0, 1)$), the system density involves
    $2^m - 1$ distinct exponential terms, and the eigenvalue ratio
    is $O(m^{-1})$.
  \item For intermediate $k$ (e.g., $k = \lceil m/2 \rceil$), the
    system density involves $O(\binom{m}{k})$ terms, and the
    eigenvalue ratio is larger.
  \item As $k \to n$ (approaching series), the number of
    system-informative terms collapses, and $\mu_{\min} \to 0$.
\end{enumerate}
\end{proposition}

The intuition is as follows.  The system density $f_T$ for a
$k$-out-of-$n$ system is a polynomial in the exponentials
$e^{-\lambda_j t}$, and the ``complexity'' of this polynomial---the
number of structurally distinct terms---determines how many independent
constraints on $\blam$ the density provides.  A richer density
(more terms) spreads information more evenly across eigenvalue
directions.

\subsection{Condition number}
\label{sec:condition-number}

\begin{definition}[Condition number of $\mathcal{I}$]
\label{def:condition-number}
The condition number of the Fisher information matrix is
\begin{equation}\label{eq:condition-number}
  \kappa(\mathcal{I}) = \frac{\mu_{\max}}{\mu_{\min}}
  = \frac{\mu_1}{\mu_m}.
\end{equation}
\end{definition}

The condition number measures the ratio of best to worst estimation
precision across linear combinations of the parameters.  When
$\kappa$ is large, the MLE is well-determined along the leading
eigenvector direction but poorly determined along the trailing
direction.  For the series system, $\kappa = \infty$
(rank-deficient); for an ideal system, $\kappa = 1$ (isotropic
information).

\begin{proposition}[Condition number determines RMSE]
\label{prop:kappa-rmse}
Under standard MLE asymptotics, the mean squared error of the MLE
$\hat{\blam}$ satisfies
\begin{equation}\label{eq:mse-trace}
  \E\bigl[\|\hat{\blam} - \blam\|^2\bigr]
  \approx \tr\bigl(\mathcal{I}^{-1}\bigr) / n
  = \frac{1}{n}\sum_{k=1}^m \frac{1}{\mu_k}.
\end{equation}
The sum is dominated by the smallest eigenvalue:
$\tr(\mathcal{I}^{-1}) \geq m / \mu_m$, so
\begin{equation}\label{eq:rmse-bound}
  \mathrm{RMSE} \geq \sqrt{m / (n\,\mu_m)}.
\end{equation}
The condition number controls the ratio of the RMSE lower bound to
the ``ideal'' RMSE (which would obtain if all eigenvalues equaled
$\mu_1$):
\begin{equation}\label{eq:rmse-ratio}
  \frac{\mathrm{RMSE}_{\text{actual}}}{\mathrm{RMSE}_{\text{ideal}}}
  \geq \sqrt{\kappa(\mathcal{I})}.
\end{equation}
\end{proposition}

\begin{proof}
By the asymptotic normality of the MLE
\citep{VanDerVaart1998, LehmannCasella1998},
$\sqrt{n}(\hat{\blam} - \blam) \xrightarrow{d} N(0, \mathcal{I}^{-1})$.
The mean squared error is
$\E[\|\hat{\blam} - \blam\|^2] \approx \tr(\mathcal{I}^{-1}) / n$.
The eigenvalue decomposition gives
$\tr(\mathcal{I}^{-1}) = \sum_k 1/\mu_k$.

If all eigenvalues equaled $\mu_1$, the trace would be $m/\mu_1$,
giving $\mathrm{RMSE}_{\text{ideal}} = \sqrt{m/(n\mu_1)}$.
The actual RMSE is at least $\sqrt{m/(n\mu_m)}$, so the ratio is
at least $\sqrt{\mu_1/\mu_m} = \sqrt{\kappa}$.
\end{proof}


%%% ====================================================================
\section{Determinant versus Condition Number}
\label{sec:det-vs-kappa}

The classical framework of optimal experimental design offers several
scalar summaries of the information matrix
\citep{Kiefer1974, Pukelsheim2006}.  In the context of system selection
for component estimation, these translate to different criteria for
choosing which system structure is ``best'' at extracting component
information.

\subsection{Optimality criteria}
\label{sec:optimality-criteria}

\begin{definition}[Design optimality criteria]
\label{def:optimality}
Given a Fisher information matrix $\mathcal{I}$ with eigenvalues
$\mu_1 \geq \cdots \geq \mu_m > 0$:
\begin{enumerate}[label=(\alph*)]
  \item \emph{D-optimality}: maximize $\det(\mathcal{I}) = \prod_k \mu_k$.
  \item \emph{A-optimality}: minimize $\tr(\mathcal{I}^{-1}) = \sum_k 1/\mu_k$.
  \item \emph{E-optimality}: maximize $\mu_m = \lambda_{\min}(\mathcal{I})$.
\end{enumerate}
\end{definition}

D-optimality corresponds to minimizing the volume of the confidence
ellipsoid; A-optimality minimizes the average variance;
E-optimality minimizes the worst-case variance.

\subsection{Why D-optimality can mislead}
\label{sec:d-misleading}

The determinant $\det(\mathcal{I}) = \prod_k \mu_k$ can be large even
when $\mu_m$ is very small, provided $\mu_1$ is sufficiently large.
A system with one enormous eigenvalue and several small ones will have
a large determinant but poor estimation quality for most parameter
directions.

\begin{proposition}[D-optimality versus RMSE]
\label{prop:d-vs-rmse}
Fix the number of components $m$ and sample size $n$.  Among
information matrices $\mathcal{I}$ with the same determinant, the
A-optimal (minimum trace of $\mathcal{I}^{-1}$) design has
$\mu_1 = \cdots = \mu_m$; that is, isotropic information.  The RMSE
ratio between the worst and best A-optimal designs with a given
determinant can be arbitrarily large.
\end{proposition}

\begin{proof}
By the AM-GM inequality,
$\tr(\mathcal{I}^{-1}) = \sum_k 1/\mu_k \geq m / (\prod_k \mu_k)^{1/m}
= m / [\det(\mathcal{I})]^{1/m}$,
with equality if and only if all eigenvalues are equal.  Conversely,
for a fixed determinant $d = \prod_k \mu_k$, the eigenvalue spectrum
$\mu_1 = d^{1/m} \cdot c^{m-1}$,
$\mu_2 = \cdots = \mu_m = d^{1/m} / c$
achieves the same determinant for any $c > 0$, but
$\tr(\mathcal{I}^{-1}) = 1/(d^{1/m} c^{m-1}) + (m-1)c / d^{1/m}$,
which diverges as $c \to \infty$.
\end{proof}

This is precisely what happens in practice.  Our numerical experiments
(Table~\ref{tab:fisher-summary}) show that the parallel(5) system has
$\det(\mathcal{I}) / \det(\mathcal{I}^{\mathrm{comp}}) = 2.27 \times 10^{-7}$,
while the 3-of-5 system has
$\det$-ratio $= 5.21 \times 10^{-10}$---three orders of magnitude
smaller.  Yet the 3-of-5 system achieves mean RMSE of $0.127$
versus $0.621$ for parallel(5).  The explanation is in the eigenvalue
spectrum: parallel(5) has $\kappa \approx 2000$, while 3-of-5 has
$\kappa \approx 269$.

\subsection{The det-ratio trap}
\label{sec:det-ratio-trap}

\begin{definition}[Determinant ratio]
\label{def:det-ratio}
The \emph{determinant ratio} (or information loss ratio) is
\begin{equation}\label{eq:det-ratio}
  \rho_D = \frac{\det(\mathcal{I}_{\text{obs}})}
               {\det(\mathcal{I}_{\text{comp}})},
\end{equation}
where $\mathcal{I}_{\text{comp}} = \diag(n/\lambda_1^2, \ldots, n/\lambda_m^2)$
is the complete-data Fisher information.
\end{definition}

The determinant ratio quantifies the total information loss from
system-level observation.  For the systems in our study, $\rho_D$
spans fourteen orders of magnitude:
from $1.46 \times 10^{-3}$ (parallel, $m = 3$) to
$4.65 \times 10^{-13}$ (consecutive-2-of-5).  But $\rho_D$ is a
poor predictor of RMSE because it conflates information along
well-estimated and poorly-estimated directions.

The condition number, by contrast, captures the \emph{shape} of the
information ellipsoid rather than its volume.  A system with a
needle-shaped ellipsoid (high $\kappa$, one long axis) can have a
larger volume than a system with a spherical ellipsoid (low $\kappa$),
yet the needle-shaped system will produce worse estimates for most
parameter directions.

\subsection{Numerical comparison}
\label{sec:numerical-comparison}

We fit the linear model
$\log(\overline{\text{RMSE}}) = \beta_0 + \beta_1 \log(\kappa)$
across the nine identifiable systems in our study (excluding the two
series systems).  The result is $R^2 = 0.91$, indicating that the
condition number alone explains 91\% of the variation in estimation
quality across system structures.  The corresponding regression using
$\log(\rho_D)$ yields $R^2 = 0.52$---substantially less predictive.

A similar regression using $\log(\mu_{\min})$ (the E-optimality
criterion) yields $R^2 = 0.86$, confirming that the smallest
eigenvalue---and hence the worst-case estimation direction---is a
better predictor than the determinant but slightly worse than the
condition number, which accounts for the scale of the leading
eigenvalue.


%%% ====================================================================
\section{Component Symmetry and Information Distribution}
\label{sec:symmetry}

\subsection{Symmetric systems have better conditioning}
\label{sec:symmetric-conditioning}

The most striking feature of our numerical experiments is the
difference in estimation quality between component-symmetric and
asymmetric systems.  Among 5-component systems, the $k$-out-of-5
structures (symmetric) uniformly outperform their consecutive-$k$-of-5
counterparts (asymmetric) in mean RMSE, despite having the same or
smaller determinant ratio.

\begin{conjecture}[Symmetric conditioning bound]
\label{conj:symmetric}
Let $\phi_S$ be a component-symmetric coherent structure function on
$m$ components and $\phi_A$ be a non-component-symmetric structure
function with the same reliability polynomial degree (i.e., the same
signature vector under i.i.d.\ components).  Then, for exponential
components with rates $\blam$ in a generic subset of $\Rset_{>0}^m$,
\begin{equation}\label{eq:symmetric-bound}
  \kappa\bigl(\mathcal{I}_S(\blam)\bigr) \leq
  \kappa\bigl(\mathcal{I}_A(\blam)\bigr).
\end{equation}
\end{conjecture}

We have verified Conjecture~\ref{conj:symmetric} numerically across
all paired comparisons in our study:

\begin{center}
\begin{tabular}{llrr}
\toprule
Symmetric & Asymmetric & $\kappa_S$ & $\kappa_A$ \\
\midrule
2-of-5 & consec-2-of-5 & $\sim 270$ & $\sim 3000$ \\
3-of-5 & consec-3-of-5 & $\sim 270$ & $\sim 4800$ \\
\bottomrule
\end{tabular}
\end{center}

The mechanism is as follows.  In a component-symmetric system, every
component plays the same structural role (same number of critical
states), so the Fisher information matrix inherits a degree of
symmetry: its off-diagonal entries are ``flattened,'' and the
eigenvalue spread is reduced.  In an asymmetric system, components at
structurally disadvantaged positions (e.g., edge components in
consecutive-$k$) have fewer critical states, contributing less to the
likelihood, and the corresponding rows and columns of $\mathcal{I}$
are depressed, widening the eigenvalue spread.

\subsection{Information deserts}
\label{sec:info-deserts}

\begin{definition}[Information desert]
\label{def:info-desert}
Component~$j$ is in an \emph{information desert} if its marginal
Fisher information $\mathcal{I}_{jj}$ satisfies
$\mathcal{I}_{jj} / \mathcal{I}_{jj}^{\mathrm{comp}} \ll 1/m$,
where $\mathcal{I}_{jj}^{\mathrm{comp}} = n/\lambda_j^2$ is the
complete-data information.  Informally, the system-level observation
provides negligible information about $\lambda_j$ relative to the
average across components.
\end{definition}

Information deserts are a signature of asymmetric systems.  In the
consecutive-2-of-5 system, component~5 (the edge component with
index at the boundary of the linear arrangement) has a per-component
RMSE of $5.746$---an order of magnitude worse than interior components
(RMSE $\approx 0.085$--$0.104$).  By contrast, in the symmetric
2-of-5 system, the RMSE spread across components is much smaller:
$0.060$ to $0.286$.

\begin{proposition}[Critical-state count and information]
\label{prop:critical-count}
The diagonal entry $\mathcal{I}_{jj}$ of the Fisher information matrix
is determined by the critical states $\mathcal{C}_j$ of component~$j$.
For a component-symmetric system, $|\mathcal{C}_j|$ is the same for
all $j$; for an asymmetric system, components with fewer critical
states have smaller $\mathcal{I}_{jj}$.
\end{proposition}

\begin{proof}
From the critical-state decomposition~\eqref{eq:f-sys-general}, the
contribution of component~$j$ to the system density is
\[
  f_T(t) \ni f_j(t) \sum_{\mathbf{x}_{-j} \in \mathcal{C}_j}
  \prod_{i \neq j} p_i(x_i, t),
\]
where $p_i(1,t) = S_i(t)$ and $p_i(0,t) = F_i(t)$.  The derivative
$\partial f_T / \partial \lambda_j$ isolates the contribution of
$f_j$ and its dependence on $\lambda_j$.  When $|\mathcal{C}_j|$ is
small, fewer state configurations contribute, and the sensitivity of
$f_T$ to $\lambda_j$ is reduced.

For a component-symmetric system,
$|\mathcal{C}_j| = |\mathcal{C}_k|$ for all $j, k$ by the symmetry
of $\phi$, so the diagonal entries of $\mathcal{I}$ are equal (at
equal rates), and the eigenvalue spread is minimized.
\end{proof}

\subsection{Path count as information proxy}
\label{sec:path-count}

The number of minimal path sets in which a component participates is a
structural proxy for its Fisher information.  A component that appears
in many minimal paths is ``structurally important''---its failure state
is constrained by many independent equations in the likelihood.

\begin{center}
\begin{tabular}{lrrrr}
\toprule
System & Paths & Cuts & Mean RMSE & Max RMSE \\
\midrule
consec-2-of-5 & 4 & 4 & 1.246 & 5.746 \\
consec-3-of-5 & 4 & 3 & 0.902 & 4.168 \\
2-of-5 & 10 & 5 & 0.128 & 0.286 \\
3-of-5 & 10 & 10 & 0.223 & 0.708 \\
bridge(5) & 4 & 4 & 0.113 & --- \\
\bottomrule
\end{tabular}
\end{center}

The $k$-out-of-5 systems, with $\binom{5}{k}$ path sets (10 for
$k = 2$ or $k = 3$), have dramatically better estimation quality than
consecutive-$k$ systems, which have only $n - k + 1 = 4$ path sets.
The bridge system, despite having only 4 path sets, achieves good RMSE
because its path/cut structure provides complementary constraints
(components appear in different roles across paths).


%%% ====================================================================
\section{Simulation Study}
\label{sec:simulations}

\subsection{Design}
\label{sec:sim-design}

We study eleven coherent system structures across three and five
components:

\begin{center}
\begin{tabular}{llrrl}
\toprule
System & Type & $m$ & Paths & Symmetric? \\
\midrule
series(3) & $k$-of-$n$ & 3 & 1 & Yes \\
parallel(3) & $k$-of-$n$ & 3 & 3 & Yes \\
2-of-3 & $k$-of-$n$ & 3 & 3 & Yes \\
series(5) & $k$-of-$n$ & 5 & 1 & Yes \\
parallel(5) & $k$-of-$n$ & 5 & 5 & Yes \\
2-of-5 & $k$-of-$n$ & 5 & 10 & Yes \\
3-of-5 & $k$-of-$n$ & 5 & 10 & Yes \\
4-of-5 & $k$-of-$n$ & 5 & 5 & Yes \\
bridge(5) & Non-SP & 5 & 4 & No \\
consec-2-of-5 & Consec. & 5 & 4 & No \\
consec-3-of-5 & Consec. & 5 & 4 & No \\
\bottomrule
\end{tabular}
\end{center}

For each structure, we generate $n = 200$ system lifetimes from
exponential components with rates
$\blam_3 = (0.2, 0.3, 0.5)$ (for $m = 3$) and
$\blam_5 = (0.1, 0.2, 0.3, 0.4, 0.5)$ (for $m = 5$), listed in
ascending order.  We fit the MLE using L-BFGS-B with Nelder--Mead
fallback and spread initialization (to break permutation symmetry),
compute the observed Fisher information matrix via numerical
differentiation, and extract eigenvalues.  Each configuration is
replicated $R = 30$ times.  Estimates are sorted in ascending order
and compared to the sorted true rates.

\subsection{Fisher information summary}
\label{sec:fisher-summary}

\begin{table}[t]
\centering
\caption{Fisher information geometry across eleven system structures.
  Rank is the numerical rank of $\hat{\mathcal{I}}$ (eigenvalue ratio
  threshold $10^{-6}$).  Det ratio is
  $\det(\hat{\mathcal{I}}) / \det(\mathcal{I}^{\mathrm{comp}})$.
  Mean RMSE is averaged over components and replications.
  $\mu_{\min}/\mu_{\max}$ is the reciprocal condition number.
  $n = 200$, $R = 30$ replications.}
\label{tab:fisher-summary}
\small
\begin{tabular}{lrrrrc}
\toprule
System & $m$ & Rank & Det ratio & Mean RMSE & $\mu_{\min}/\mu_{\max}$ \\
\midrule
series(3) & 3 & 1 & $1.01 \times 10^{-24}$ & $0.086^\dagger$ & $\approx 0$ \\
parallel(3) & 3 & 3 & $1.46 \times 10^{-3}$ & 0.111 & 0.0056 \\
2-of-3 & 3 & 3 & $6.33 \times 10^{-5}$ & 0.104 & 0.0067 \\
\midrule
series(5) & 5 & 1 & $1.09 \times 10^{-50}$ & $0.063^\dagger$ & $\approx 0$ \\
parallel(5) & 5 & 5 & $2.27 \times 10^{-7}$ & 0.621 & 0.0005 \\
2-of-5 & 5 & 5 & $1.09 \times 10^{-8}$ & 1.421 & 0.0020 \\
3-of-5 & 5 & 5 & $5.21 \times 10^{-10}$ & 0.127 & 0.0037 \\
4-of-5 & 5 & 5 & $1.71 \times 10^{-12}$ & 0.140 & 0.0037 \\
bridge(5) & 5 & 5 & $3.82 \times 10^{-11}$ & 0.113 & 0.0012 \\
consec-2-of-5 & 5 & 5 & $4.65 \times 10^{-13}$ & 3.088 & 0.0004 \\
consec-3-of-5 & 5 & 5 & $4.39 \times 10^{-10}$ & 1.798 & 0.0004 \\
\bottomrule
\multicolumn{6}{l}{\footnotesize $^\dagger$ Series systems: non-identifiable; RMSE for total rate only.}
\end{tabular}
\end{table}

Table~\ref{tab:fisher-summary} reveals several key patterns:
\begin{enumerate}[leftmargin=2em]
  \item \emph{Rank}:  Series systems have rank 1 (confirming
    Theorem~\ref{thm:series-rank}); all other systems have full rank
    (confirming Theorem~\ref{thm:kofn-rank} and
    Corollary~\ref{cor:general-coherent}).

  \item \emph{Det ratio is not monotone with RMSE}:
    The 3-of-5 system has a smaller det ratio
    ($5.21 \times 10^{-10}$) than parallel(5)
    ($2.27 \times 10^{-7}$), yet achieves much better RMSE
    (0.127 vs.\ 0.621).

  \item \emph{Condition number tracks RMSE}:
    The reciprocal condition number $\mu_{\min}/\mu_{\max}$ is
    strongly correlated with estimation quality: the best systems
    (3-of-5, 4-of-5) have the largest $\mu_{\min}/\mu_{\max}$,
    while the worst (consecutive-$k$) have the smallest.

  \item \emph{Asymmetric systems are outliers}:
    The consecutive-$k$ systems have both the smallest
    $\mu_{\min}/\mu_{\max}$ and the largest RMSE, driven by
    information deserts at edge components.
\end{enumerate}

\subsection{Eigenvalue profiles}
\label{sec:ev-profiles}

The eigenvalue profiles for three representative 5-component systems
illustrate the key differences:

\begin{center}
\begin{tabular}{lrrrrrl}
\toprule
System & $\mu_1$ & $\mu_2$ & $\mu_3$ & $\mu_4$ & $\mu_5$ & RMSE \\
\midrule
3-of-5 & 1911 & 16.7 & 8.2 & 7.8 & 7.1 & 0.127 \\
parallel(5) & 18022 & 444 & 29 & 13 & 9 & 0.621 \\
consec-2-of-5 & 1222 & 40.5 & 5.1 & 1.0 & 0.4 & 3.088 \\
\bottomrule
\end{tabular}
\end{center}

In the 3-of-5 system, the trailing eigenvalues ($\mu_2$ through
$\mu_5$) are relatively close together and not too far below $\mu_1$
(condition number $\approx 269$).  The parallel(5) system has a much
larger $\mu_1$ but steeper decay in the trailing eigenvalues
($\kappa \approx 2000$).  The consecutive-2-of-5 system has the
steepest decay ($\kappa \approx 3000$), with $\mu_5 = 0.4$---three
orders of magnitude below $\mu_1$.

The leading eigenvalue $\mu_1$ is always by far the largest, confirming
Proposition~\ref{prop:leading-ev}: the total rate is always
well-determined.  The practical estimation difficulty is entirely in
the trailing eigenvalues.

\subsection{Per-component estimation quality}
\label{sec:per-component}

For asymmetric systems, the estimation quality varies dramatically
across components.  Table~\ref{tab:per-component} compares the
per-component RMSE for consecutive-2-of-5 (asymmetric) and 2-of-5
(symmetric).

\begin{table}[t]
\centering
\caption{Per-component RMSE for symmetric (2-of-5) and asymmetric
  (consec-2-of-5) systems.  True rates are
  $\blam = (0.1, 0.2, 0.3, 0.4, 0.5)$ sorted ascending.
  Component indices correspond to sorted positions.}
\label{tab:per-component}
\small
\begin{tabular}{lrrrrr}
\toprule
 & $\hat{\lambda}_1$ & $\hat{\lambda}_2$ & $\hat{\lambda}_3$
 & $\hat{\lambda}_4$ & $\hat{\lambda}_5$ \\
\midrule
2-of-5 & 0.081 & 0.060 & 0.096 & 0.117 & 0.286 \\
consec-2-of-5 & 0.085 & 0.086 & 0.104 & 0.211 & 5.746 \\
\bottomrule
\end{tabular}
\end{table}

The consecutive-2-of-5 system exhibits a catastrophic RMSE of $5.746$
for the largest estimated rate (corresponding to the fastest-failing
component at the edge of the linear arrangement).  This is an
information desert: the edge component participates in only one
minimal cut set, so the likelihood has minimal sensitivity to its rate.
In the symmetric 2-of-5 system, the worst per-component RMSE is
$0.286$---still elevated for the fastest component, but not
catastrophically so.

\subsection{Condition number as predictor}
\label{sec:kappa-predictor}

\begin{figure}[t]
\centering
% [Figure placeholder: log-log plot of mean RMSE vs condition number]
% x-axis: log10(kappa), y-axis: log10(mean RMSE)
% 9 points (excluding series systems), linear fit with R^2 annotation
\fbox{\parbox{0.8\textwidth}{\centering
  \vspace{3em}
  {\small Log-log regression of mean RMSE on condition number $\kappa$
  across nine identifiable structures.\\ Fitted line:
  $\log_{10}(\text{RMSE}) = \beta_0 + \beta_1 \log_{10}(\kappa)$,
  $R^2 = 0.91$.}
  \vspace{3em}
}}
\caption{Mean RMSE versus condition number of the observed Fisher
  information matrix, log-log scale, across nine identifiable
  structures.  The strong linear relationship ($R^2 = 0.91$) confirms
  that the condition number is a near-complete summary of estimation
  difficulty.}
\label{fig:kappa-rmse}
\end{figure}

Figure~\ref{fig:kappa-rmse} displays the log-log relationship between
the condition number and mean RMSE across the nine identifiable
structures.  The fitted regression has $R^2 = 0.91$, confirming that
the condition number is a near-complete summary of estimation
difficulty.  By comparison, the determinant ratio yields
$R^2 = 0.52$.


%%% ====================================================================
\section{Discussion}
\label{sec:discussion}

\subsection{Summary of findings}

This paper establishes that the eigenvalue structure of the Fisher
information matrix provides a complete characterization of the
difficulty of component rate estimation from system-level lifetime
data.  The key insights are:

\begin{enumerate}[leftmargin=2em]
  \item \emph{Full rank is necessary but not sufficient.}  All
    non-series coherent systems with distinct exponential rates have
    full-rank Fisher information (Theorem~\ref{thm:kofn-rank}), but
    the condition number $\kappa$ varies by more than an order of
    magnitude across system structures, translating directly to
    estimation quality.

  \item \emph{The leading eigenvalue is universal.}  Across all
    structures, the leading eigenvalue corresponds to the total-rate
    direction $\bone/\sqrt{m}$ and is always large
    (Proposition~\ref{prop:leading-ev}).  The total rate
    $\lambda_+ = \sum_j \lambda_j$ is always well-estimated; the
    difficulty lies in resolving individual rates.

  \item \emph{D-optimality is misleading.}  The determinant of the
    Fisher information can be large when one eigenvalue dominates,
    making the system appear more informative than it is for practical
    estimation.  The condition number is a far better predictor of RMSE
    ($R^2 = 0.91$ vs.\ $0.52$ for the determinant ratio).

  \item \emph{Symmetry helps.}  Component-symmetric systems
    ($k$-out-of-$n$) have uniformly better conditioning than asymmetric
    systems (consecutive-$k$) with comparable complexity, because
    symmetry distributes information evenly across components.

  \item \emph{Asymmetry creates information deserts.}  In
    consecutive-$k$ systems, edge components can have RMSE values
    orders of magnitude worse than interior components, making them
    effectively unestimable at practical sample sizes.
\end{enumerate}

\subsection{Connection to optimal design theory}
\label{sec:design-connection}

The parallel between our problem and classical optimal experimental
design \citep{Pukelsheim2006} is informative.  In experimental design,
the analyst chooses a design (allocation of observations to
experimental conditions) to optimize a criterion applied to the
Fisher information matrix.  In our setting, the ``design'' is the
system structure, which is typically fixed by engineering requirements
rather than chosen by the statistician.  Nonetheless, the design
theory framework provides useful vocabulary: our finding that
D-optimality misleads and E-optimality is more predictive of RMSE
parallels known results in experimental design
\citep{Atkinson2007}, where E-optimal designs are preferred when the
goal is to estimate all parameters well, not just to minimize the
volume of the confidence region.

One actionable implication: when an engineer has the freedom to choose
among redundancy structures (e.g., $k$-out-of-$n$ with different $k$),
our results suggest choosing the structure with the best-conditioned
Fisher information matrix---typically an intermediate $k$---rather
than the one with the largest determinant.

\subsection{The role of the signature}
\label{sec:role-of-signature}

The system signature $\mathbf{s}$ \citep{Samaniego1985} encodes where
in the order-statistic sequence the system fails.  For i.i.d.\
components, signatures fully characterize system reliability.  For
heterogeneous components, the signature does not directly apply, but
it provides intuition: a system whose signature concentrates near the
middle of the failure sequence (e.g., 3-of-5 with
$\mathbf{s} = (0, 0, 1, 0, 0)$) tends to have better conditioning
than one concentrated at the extremes (parallel with
$\mathbf{s} = (0, 0, 0, 0, 1)$ or series with
$\mathbf{s} = (1, 0, 0, 0, 0)$).

This is because a ``middle'' failure position creates maximum
mixing of survival and failure probabilities in the system density,
generating the richest set of constraints on the component rates.

\subsection{Extensions}
\label{sec:extensions}

Several directions merit further investigation:

\paragraph{Beyond exponential components.}
The exponential assumption is restrictive.  For Weibull components
with shape and scale parameters, the Fisher information matrix is
$2m \times 2m$, and the eigenvalue structure is richer.  Preliminary
numerical experiments (not reported here) suggest that the condition
number remains a strong predictor of RMSE, but the relationship
between system structure and conditioning changes with the shape
parameter.

\paragraph{Dependent components.}
Under component dependence (e.g., via copulas or frailty models),
the critical-state decomposition~\eqref{eq:f-sys-general} no longer
applies directly.  The Fisher information geometry in this setting is
an open problem.

\paragraph{Information-optimal system selection.}
Given a fixed number of components $m$ and a cost constraint, which
coherent structure maximizes the minimum eigenvalue of $\mathcal{I}$?
This is an information-theoretic version of the system design problem.
Our numerical results suggest that $k$-out-of-$n$ systems with
$k \approx n/2$ are close to E-optimal, but a formal proof is lacking.

\paragraph{Sample size requirements.}
The condition number $\kappa$ determines how much data is needed for
accurate estimation.  A practical guideline would be:
$n \geq c \cdot \kappa \cdot m$, where $c$ is a constant depending
on the desired precision.  Making this precise requires a refined
finite-sample analysis.

\paragraph{Censored system observations.}
When the system lifetime itself is right-censored (as in life testing
with a fixed observation window), the Fisher information changes.  The
interaction between system-level censoring and the eigenvalue structure
of $\mathcal{I}$ is an interesting question for future work.

\subsection{Open problems}
\label{sec:open-problems}

\begin{enumerate}[leftmargin=2em]
  \item \emph{Prove Conjecture~\ref{conj:symmetric}}: does component
    symmetry always improve conditioning?  The conjecture is supported
    by all our numerical experiments but lacks a proof.

  \item \emph{Closed-form eigenvalue expressions}: for $k$-out-of-$n$
    with i.i.d.\ exponential components, can the eigenvalues of
    $\mathcal{I}$ be expressed in terms of the harmonic numbers
    $H_m = \sum_{j=1}^m 1/j$?

  \item \emph{Minimax-optimal structure}: among all coherent structures
    on $m$ components, which minimizes $\kappa(\mathcal{I})$?

  \item \emph{Eigenvalue structure for non-exponential families}:
    how does the relationship between system structure and conditioning
    change for Weibull, gamma, or log-normal components?
\end{enumerate}


%%% ====================================================================
\section*{Acknowledgments}

Computational experiments were performed using R.  The coherent system
engine and general likelihood code are available at
\url{https://github.com/queelius}.


%%% ====================================================================
\appendix
\section{Supplementary Proofs}
\label{sec:appendix}

\subsection{Proof of Proposition~\ref{prop:leading-ev}}
\label{sec:proof-leading-ev}

We provide a more detailed argument for the alignment of the leading
eigenvector with the total-rate direction.

Consider the score function for a single observation $T = t$:
\[
  s_j(t; \blam) = \frac{\partial}{\partial \lambda_j} \log f_T(t; \blam)
  = \frac{1}{f_T(t; \blam)} \cdot \frac{\partial f_T(t; \blam)}{\partial \lambda_j}.
\]
The Fisher information matrix is $\mathcal{I}_{jk} = \E[s_j(T) s_k(T)]$.

For any coherent system with exponential components, the system density
$f_T(t; \blam)$ is a sum of terms of the form
$c_{S} \, \lambda_j \, e^{-\alpha_S t}$, where $\alpha_S$ is a sum
of rates involving the rates of all components.  The partial derivative
$\partial f_T / \partial \lambda_j$ therefore produces terms that
involve $\lambda_j$ and rate sums $\alpha_S$ that include $\lambda_j$.

The key observation is that for large $t$, the dominant terms in
$f_T(t)$ and its derivatives involve the smallest possible rate sum
$\alpha_S$, which is bounded below by $\lambda_1$ (the smallest
component rate).  This dominant behavior is shared across all partial
derivatives, creating a strong positive correlation among the score
components $s_1, \ldots, s_m$.  This correlation manifests as a large
rank-1 component $\alpha \bone \bone^\top$ in $\mathcal{I}$.

More precisely, decompose
\[
  s_j(t) = g(t) + h_j(t),
\]
where $g(t)$ is the common component (arising from the dominant
exponential terms) and $h_j(t)$ captures the $j$-specific deviation.
Then
\[
  \mathcal{I}_{jk} = \E[g^2] + \E[g(h_j + h_k)] + \E[h_j h_k].
\]
The first term is $\E[g^2] \cdot \bone\bone^\top$ (rank 1), and the
remaining terms form the residual $\mathbf{R}$.  Since $\E[g^2]$
captures the dominant exponential behavior, it is typically much
larger than the entries of $\mathbf{R}$, giving $\mathbf{v}_1 \approx
\bone / \sqrt{m}$.

\subsection{Verification of eigenvalue profiles}
\label{sec:verify-ev}

The eigenvalue profiles reported in
\S\ref{sec:ev-profiles} were computed from the observed Fisher
information matrix $\hat{\mathcal{I}} = -\nabla^2 \ell(\hat{\blam})$
using numerical second derivatives via the Richardson extrapolation
method implemented in the R package \texttt{numDeriv}.  Eigenvalues
were computed using the spectral decomposition with symmetry enforced
($\hat{\mathcal{I}}$ is symmetric by construction).

To guard against numerical artifacts, we verified that:
\begin{enumerate}[label=(\roman*)]
  \item all eigenvalues are non-negative (as required for a positive
    semidefinite information matrix);
  \item the rank determination (eigenvalue ratio threshold $10^{-6}$)
    is robust to perturbations of the threshold;
  \item the eigenvalue profiles are consistent across Monte Carlo
    replications (coefficient of variation $< 15\%$ for all nonzero
    eigenvalues).
\end{enumerate}


\bibliographystyle{plainnat}
\bibliography{references}

\end{document}
